\documentclass[11pt]{article}

\usepackage[margin=1in]{geometry}
\usepackage{amsmath,amssymb,amsthm,mathtools}
\usepackage{bm}
\usepackage{hyperref}

\title{Scaling Relation Between $G$ and $G_\alpha$ for the 1D Periodic Dirichlet--Neumann Operator}
\author{}
\date{}

\newtheorem{proposition}{Proposition}
\newtheorem{remark}{Remark}

\begin{document}
\maketitle

\section{Original problem on $[0,2\pi]$}

Let $x \in [0,2\pi]$ be the horizontal variable, with periodic boundary conditions, and let the fluid domain be
\[
\Omega_\eta
=
\{(x,y): x \in [0,2\pi],\ -h < y < \eta(x)\}.
\]
Given a surface profile $\eta(x)$ and boundary data $\xi(x)$, let $\phi(x,y)$ solve
\begin{equation}
\begin{cases}
\phi_{xx} + \phi_{yy} = 0, & (x,y) \in \Omega_\eta,\\[0.3em]
\phi_y(x,-h)=0,\\[0.3em]
\phi(x,\eta(x))=\xi(x),
\end{cases}
\label{eq:orig_bvp}
\end{equation}
with $2\pi$-periodicity in $x$.

The Dirichlet--Neumann operator is defined by
\begin{equation}
G(\eta)\xi
=
\left( -\eta_x,\, 1 \right)\cdot \nabla \phi \Big|_{y=\eta(x)}
=
\phi_y(x,\eta(x))-\eta_x(x)\phi_x(x,\eta(x)).
\label{eq:G_def}
\end{equation}

\section{Rescaled problem on $[0,2\pi\alpha]$}

Now fix $\alpha>0$ and consider the new horizontal domain
\[
x \in [0,2\pi\alpha].
\]
We also scale the depth by the same factor, so the new fluid domain is
\[
\Omega^{(\alpha)}_\eta
=
\{(x,y): x \in [0,2\pi\alpha],\ -h\alpha < y < \eta(x)\}.
\]
Given a surface profile $\eta(x)$ and boundary data $\xi(x)$ on $[0,2\pi\alpha]$, let $\Phi(x,y)$ solve
\begin{equation}
\begin{cases}
\Phi_{xx} + \Phi_{yy} = 0, & (x,y) \in \Omega^{(\alpha)}_\eta,\\[0.3em]
\Phi_y(x,-h\alpha)=0,\\[0.3em]
\Phi(x,\eta(x))=\xi(x),
\end{cases}
\label{eq:alpha_bvp}
\end{equation}
with $2\pi\alpha$-periodicity in $x$.

We define the corresponding Dirichlet--Neumann operator by
\begin{equation}
G_\alpha(\eta)\xi
=
\left( -\eta_x,\, 1 \right)\cdot \nabla \Phi \Big|_{y=\eta(x)}
=
\Phi_y(x,\eta(x))-\eta_x(x)\Phi_x(x,\eta(x)).
\label{eq:Galpha_def}
\end{equation}

\section{Change of variables}

To compare $G$ and $G_\alpha$, introduce the rescaled variables
\begin{equation}
x = \alpha X,
\qquad
y = \alpha Y,
\label{eq:change_var}
\end{equation}
so that
\[
X \in [0,2\pi].
\]
Given a function $\Phi(x,y)$ on the rescaled domain, define
\begin{equation}
\phi(X,Y) := \Phi(\alpha X,\alpha Y).
\label{eq:phi_pullback}
\end{equation}
We also define the pulled-back surface and Dirichlet data by
\begin{equation}
\widetilde{\eta}(X) := \frac{1}{\alpha}\eta(\alpha X),
\qquad
\widetilde{\xi}(X) := \xi(\alpha X).
\label{eq:eta_xi_tilde}
\end{equation}

\subsection*{Step 1: the domain becomes $[0,2\pi] \times [-h,\widetilde{\eta}(X)]$}

Indeed, the bottom boundary $y=-h\alpha$ becomes
\[
Y=-h,
\]
and the free surface $y=\eta(x)$ becomes
\[
\alpha Y = \eta(\alpha X)
\quad \Longleftrightarrow \quad
Y = \frac{1}{\alpha}\eta(\alpha X)=\widetilde{\eta}(X).
\]
Hence the transformed domain is
\[
\widetilde{\Omega}
=
\{(X,Y): X\in[0,2\pi],\ -h<Y<\widetilde{\eta}(X)\}.
\]

\subsection*{Step 2: the Laplace equation is preserved}

From \eqref{eq:phi_pullback},
\[
\phi_X(X,Y) = \alpha\, \Phi_x(\alpha X,\alpha Y),
\qquad
\phi_Y(X,Y) = \alpha\, \Phi_y(\alpha X,\alpha Y).
\]
Differentiating once more,
\[
\phi_{XX}(X,Y)=\alpha^2 \Phi_{xx}(\alpha X,\alpha Y),
\qquad
\phi_{YY}(X,Y)=\alpha^2 \Phi_{yy}(\alpha X,\alpha Y).
\]
Therefore
\[
\phi_{XX}+\phi_{YY}
=
\alpha^2(\Phi_{xx}+\Phi_{yy})=0.
\]
So $\phi$ is harmonic in $\widetilde{\Omega}$.

\subsection*{Step 3: the bottom boundary condition is preserved}

Since
\[
\phi_Y(X,Y)=\alpha\,\Phi_y(\alpha X,\alpha Y),
\]
at $Y=-h$ we get
\[
\phi_Y(X,-h)=\alpha\,\Phi_y(\alpha X,-h\alpha)=0.
\]

\subsection*{Step 4: the Dirichlet boundary condition becomes $\phi(X,\widetilde{\eta}(X))=\widetilde{\xi}(X)$}

On the free surface,
\[
\phi(X,\widetilde{\eta}(X))
=
\Phi(\alpha X,\alpha \widetilde{\eta}(X))
=
\Phi(\alpha X,\eta(\alpha X))
=
\xi(\alpha X)
=
\widetilde{\xi}(X).
\]

Therefore $\phi$ solves the boundary value problem
\begin{equation}
\begin{cases}
\phi_{XX}+\phi_{YY}=0, & -h<Y<\widetilde{\eta}(X),\\[0.3em]
\phi_Y(X,-h)=0,\\[0.3em]
\phi(X,\widetilde{\eta}(X))=\widetilde{\xi}(X),
\end{cases}
\label{eq:transformed_bvp}
\end{equation}
with $2\pi$-periodicity in $X$.

By definition, this means that $\phi$ is exactly the harmonic extension associated with the operator
\[
G(\widetilde{\eta})\widetilde{\xi}.
\]

\section{Transformation of the Dirichlet--Neumann operator}

We now compute the relation between $G_\alpha$ and $G$.

From \eqref{eq:phi_pullback},
\[
\Phi_x(\alpha X,\alpha Y)=\frac{1}{\alpha}\phi_X(X,Y),
\qquad
\Phi_y(\alpha X,\alpha Y)=\frac{1}{\alpha}\phi_Y(X,Y).
\]
Also, from \eqref{eq:eta_xi_tilde},
\[
\widetilde{\eta}_X(X)
=
\frac{d}{dX}\left( \frac{1}{\alpha}\eta(\alpha X)\right)
=
\eta_x(\alpha X).
\]
Hence
\[
\eta_x(\alpha X)=\widetilde{\eta}_X(X).
\]

Evaluating \eqref{eq:Galpha_def} at the point $x=\alpha X$, we obtain
\begin{align*}
\bigl(G_\alpha(\eta)\xi\bigr)(\alpha X)
&=
\Phi_y(\alpha X,\eta(\alpha X))
-
\eta_x(\alpha X)\Phi_x(\alpha X,\eta(\alpha X))\\
&=
\frac{1}{\alpha}\phi_Y(X,\widetilde{\eta}(X))
-
\widetilde{\eta}_X(X)\frac{1}{\alpha}\phi_X(X,\widetilde{\eta}(X))\\
&=
\frac{1}{\alpha}
\Bigl(
\phi_Y(X,\widetilde{\eta}(X))
-
\widetilde{\eta}_X(X)\phi_X(X,\widetilde{\eta}(X))
\Bigr).
\end{align*}
Using the definition of $G$ in \eqref{eq:G_def}, this becomes
\begin{equation}
\boxed{
\bigl(G_\alpha(\eta)\xi\bigr)(\alpha X)
=
\frac{1}{\alpha}
\bigl(G(\widetilde{\eta})\widetilde{\xi}\bigr)(X),
}
\label{eq:main_relation_pointwise}
\end{equation}
where
\[
\widetilde{\eta}(X)=\frac{1}{\alpha}\eta(\alpha X),
\qquad
\widetilde{\xi}(X)=\xi(\alpha X).
\]

\section{Operator form}

Define the pullback operator $T_\alpha$ by
\[
(T_\alpha f)(X):=f(\alpha X),
\]
which maps $2\pi\alpha$-periodic functions to $2\pi$-periodic functions.

Then
\[
\widetilde{\xi}=T_\alpha \xi,
\qquad
\widetilde{\eta}=\frac{1}{\alpha}T_\alpha \eta.
\]
Equation \eqref{eq:main_relation_pointwise} can be written as
\begin{equation}
\boxed{
T_\alpha\bigl(G_\alpha(\eta)\xi\bigr)
=
\frac{1}{\alpha}
G\!\left( \frac{1}{\alpha}T_\alpha \eta \right)(T_\alpha \xi).
}
\label{eq:main_relation_operator}
\end{equation}
Equivalently,
\begin{equation}
\boxed{
G_\alpha(\eta)
=
T_\alpha^{-1}
\left[
\frac{1}{\alpha}
G\!\left( \frac{1}{\alpha}T_\alpha \eta \right)
\right]
T_\alpha.
}
\label{eq:main_relation_conjugation}
\end{equation}

\section{Flat-surface check}

As a consistency check, let $\eta\equiv 0$. Then on the $2\pi$-periodic strip,
\[
G(0)=|D|\tanh(h|D|),
\qquad D=-i\partial_X.
\]
On the $2\pi\alpha$-periodic strip with depth $h\alpha$,
\[
G_\alpha(0)=|D_x|\tanh(h\alpha |D_x|),
\qquad D_x=-i\partial_x.
\]
Under the change $x=\alpha X$, we have
\[
D_x=\frac{1}{\alpha}D.
\]
Therefore
\[
|D_x|\tanh(h\alpha |D_x|)
=
\frac{1}{\alpha}|D|\tanh(h|D|),
\]
which is exactly consistent with \eqref{eq:main_relation_operator}.

\begin{proposition}
Let $G$ be the Dirichlet--Neumann operator on the $2\pi$-periodic domain with depth $h$, and let $G_\alpha$ be the Dirichlet--Neumann operator on the $2\pi\alpha$-periodic domain with depth $h\alpha$. Then, with
\[
(T_\alpha f)(X)=f(\alpha X),
\]
one has
\[
T_\alpha\bigl(G_\alpha(\eta)\xi\bigr)
=
\frac{1}{\alpha}
G\!\left( \frac{1}{\alpha}T_\alpha\eta \right)(T_\alpha\xi).
\]
Equivalently,
\[
G_\alpha(\eta)
=
T_\alpha^{-1}
\left[
\frac{1}{\alpha}
G\!\left( \frac{1}{\alpha}T_\alpha\eta \right)
\right]
T_\alpha.
\]
\end{proposition}

\begin{remark}
The key point is that both the horizontal variable and the vertical variable are scaled by the same factor $\alpha$. Because of this, the transformed boundary value problem again has depth $h$ on the reference interval $[0,2\pi]$, and the Dirichlet--Neumann operator picks up exactly one prefactor $1/\alpha$.
\end{remark}

\section{What to do next}
We will now start to investigate the case when the width and depth of the domain are scaled differently.

\end{document}